% 1-WL 局限：六环 C6 与两个不相交三角形；初始同标记 A，每点每轮看到两个同类邻居，签名恒为 (A,{{A,A}})。
\begin{tikzpicture}[panelmark/.style={font=\figurefont\fontsize{8}{10}\selectfont,text=FigureInk,anchor=north,inner sep=0pt,outer sep=0pt},
font=\figurefont\fontsize{8}{10}\selectfont,
  vtx/.style={circle,draw=FigureStroke,line width=.8pt,minimum size=6.2mm,inner sep=0pt,text=FigureInk,fill=FigureBlueFill},
  edge/.style={draw=FigureStroke,line width=.85pt},
  note/.style={font=\FigureNoteFont\fontsize{7.5}{9}\selectfont,text=FigureMuted,align=center}]
\begin{scope}[local bounding box=panelrow]
\begin{scope}[local bounding box=panela]
  \foreach \i in {1,...,6} \node[vtx] (a\i) at ({90-60*(\i-1)}:1.05) {$A$};
  \foreach \i/\j in {1/2,2/3,3/4,4/5,5/6,6/1} \draw[edge] (a\i)--(a\j);

\end{scope}
\begin{scope}[xshift=3.3cm,local bounding box=panelb]
  \foreach \x/\p in {0/b,2.25/c}{
    \node[vtx] (\p1) at (\x,0.75) {$A$};
    \node[vtx] (\p2) at (\x-0.65,-0.4) {$A$};
    \node[vtx] (\p3) at (\x+0.65,-0.4) {$A$};
    \draw[edge] (\p1)--(\p2)--(\p3)--(\p1);
  }

\end{scope}
\end{scope}
% Subfigure marks: 8 pt, centred below each panel; common top 3 mm below this row.
\coordinate (panelrowbase) at ($(panelrow.south)+(0,-3mm)$);
\node[panelmark] at (panela.center |- panelrowbase) {(a)};
\node[panelmark] at (panelb.center |- panelrowbase) {(b)};
\end{tikzpicture}
